\begin{lemma}[Relabelling triplets and their matching intersections]
Let $e:V\to V'$ be a bijection between finite vertex types.
Let $F$ be a graph on $V$ with disjoint finite parts $U,K$.
Put $U'=e(U)$, $K'=e(K)$ and $F'=e(F)$, where an edge of $e(F)$
means that its inverse images are adjacent in $F$.
There are bijections $e_U:U\to U'$, $e_K:K\to K'$ agreeing with $e$
on underlying vertices, disjointness $U'\cap K'=\varnothing$, and a
bijection $e_M$ between the perfect matchings of $(F;U,K)$ and
$(F';U',K')$, with
\[
 e_M(\sigma)(e_U(u))=e_K(\sigma(u)).
\]
For every error and density, $(F;U,K)$ is super-regular if and only
if $(F';U',K')$ is super-regular with those same parameters.
For any two further graphs $F_{13},F_{23}$, part $Z$, size $m$, and
parameters, the near-equal regular triplet on $(U,K,Z)$ is equivalent
to the corresponding triplet of relabelled graphs on $(U',K',e(Z))$.
For every perfect matching $\sigma$, error $\delta$, and density $p$,
its matching-induced intersection is $(\delta,p)$-super-regular on
the full canonical subtype copies of $U,Z$ if and only if the
intersection induced by $e_M(\sigma)$ in the relabelled graphs is
$(\delta,p)$-super-regular on the full subtype copies of $U',e(Z)$.
\end{lemma}
\begin{proof}
Restriction of $e$ to any finite part is injective and surjective onto
its image, yielding $e_U,e_K$ and likewise $e_Z:Z\to e(Z)$.
If $e(u)=e(k)$ with $u\in U$, $k\in K$, injectivity would give
$u=k$, contrary to disjointness; thus the new parts are disjoint.
Define $e_M(\sigma)=e_K\circ\sigma\circ e_U^{-1}$.
This is bijective because all three factors are bijective.
Its edges are edges of $F'$ exactly when the original matching edges
are edges of $F$, by the definition of the relabelled graph.
Conversely $\sigma'=e_K^{-1}\circ\tau\circ e_U$ is a perfect matching
whenever $\tau$ is one in $F'$. The two constructions are inverse,
proving the matching bijection and the displayed evaluation identity.
Here the meaning of a perfect matching is \noderef{PerfectMatching}.

For any rectangle $P\subseteq U$, $Q\subseteq K$, the map
$(u,k)\mapsto(e(u),e(k))$ is a bijection from its edges to the edges
of the corresponding rectangle in $F'$. Also $|e(P)|=|P|$ and
$|e(Q)|=|Q|$. Consequently the edge densities, defined by
\noderef{PairDensity}, are equal, including empty rectangles.
Every rectangle in the image parts arises in this way by inverse image.
For each vertex the same map gives a bijection of its neighbours
in the opposite part. Thus degrees, part sizes, and the rectangle
size tests in \noderef{SuperRegularPair} are unchanged. All its
inequalities transfer in both directions. Apply this argument also
to the two other pairs. Cardinalities of the three parts and their
disjointness are unchanged, so the three possible near-equal size
relations are unchanged as well. For an edge of the current pair,
$z\mapsto e(z)$ is a bijection between the tested common
neighbourhoods in $Z$ and $e(Z)$. Hence their cardinalities and the
window normalized by $|Z|$ are identical. These facts verify every
clause of \noderef{NearEqualRegularTriplet} in both directions.

Finally use the bijection $e_U\oplus e_Z$ between the two subtype
sum types. By \noderef{MatchingInducedGraph}, a cross edge in the
first intersection means
$F_{13}(u,z)$ and $F_{23}(\sigma(u),z)$.
In the second intersection these conditions are precisely
$e(F_{13})(e_U(u),e_Z(z))$ and
$e(F_{23})(e_M(\sigma)(e_U(u)),e_Z(z))$, by the evaluation identity.
Thus this sum bijection preserves and reflects every cross edge and
maps each full copied part exactly onto its counterpart. The same
degree and rectangle counting argument proves the last equivalence.
No probability statement is needed for this relabelling result.
\end{proof}
